% !TEX encoding = UTF-8 Unicode
% Signature: Adriano Fagiolini, Template_Exams@UNIPA, 2021-Dec-21
% Modified 2023-Jan-03 to Project form

\documentclass[12pt]{exam} 

\def\theCourse{Controlli Automatici}
\def\theExamAcronym{CA}

% !TEX encoding = UTF-8 Unicode
\usepackage[a4paper, top=25mm, bottom=25mm, left=25mm, right=25mm]{geometry}
\usepackage{microtype} % Tweak font spacing for aesthetics
%\usepackage{cmbright} % Font
\usepackage[utf8]{inputenc}
\usepackage[T1]{fontenc} 
\usepackage[italian]{babel} 
\usepackage[colorlinks=true, allcolors=black]{hyperref}
\usepackage{subfigure} % Multiple figures environment
\usepackage{amsmath,amssymb}
\usepackage{graphicx} % Include graphics
\usepackage{booktabs} % Professional tables
\usepackage[dvipsnames]{xcolor}  % Advanced colours
\usepackage{soul} % Highlight text and other text things
\usepackage{caption} % Manage captions in floats
\usepackage{textcomp} % Includes euro symbol
\usepackage{ifthen} % Use of conditional \if
\usepackage{fmtcount} % Convert integers to words
\usepackage[per-mode = symbol]{siunitx}  % To use SI units
\usepackage[version=4]{mhchem}  % For writing chemical symbols
\usepackage{totcount} % Display points counter at beginning 
\usepackage{xstring} % To manipulate strings
\usepackage{multicol} % Multiple column environment (FS)
%\usepackage{lipsum,turnthepage}
\usepackage{wasysym} % Symbols for multiple choice questions
\usepackage{fontawesome}
\usepackage{fontawesome}
\usepackage[most]{tcolorbox} % to circle the text of the email
\usepackage{multicol}

% DOCUMENT DEFINITIONS
\pointsinrightmargin  % Points on the right margin
\pointformat{\bfseries(\themarginpoints)}  % Points in bold and bracket style
\renewcommand{\partlabel}{\bfseries{(\thepartno)}} % Bold part numbers
\renewcommand{\subpartlabel}{\bfseries{\thesubpart)}} % Bold part numbers
\pagestyle{headandfoot} % Solution header and footer
\header{}{}{} % Solution header
\footer{}{}{\theCourse{} -- \theExamDate{} -- Gruppo \theGroupID{}} 
%\graphicspath{{img/}} % Directory where pictures are stored
\definecolor{myRed}{RGB}{184,0,1} % Define new colour
\definecolor{myGrey}{RGB}{211,211,211}  % Define new colour
\definecolor{verylightgray}{RGB}{241,241,241}  % Define new colour
\checkedchar{\large $\CIRCLE$}
\setlength{\fboxrule}{0pt}

% VERSION TRACKING
\newcommand{\theVersion}{1.1}
\newcommand{\theVersionDate}{13-Dec-2021}

% COURSE
\def\theExamEmail{esame.mirpalab@gmail.com}
\newcommand{\theCourseCode}{CODE12345}
\newcommand{\theDiet}{Diet} % EDIT - Exam diet (Main OR Resit)
\newcommand{\examOC}{Open}
\newcommand{\useCalculators}{Yes}
\newcommand{\formulaSheet}{Yes}
\newcommand{\hasSections}{Yes}
\newcommand{\separateBooks}{No} % Answer sections in separate books
\newcommand{\numberQuestions}{All}
\newcommand{\withSolutions}{Yes}

% NEW COMMANDS
\newtotcounter{pointsCounter} % Counter for total number of points in parts
\newtotcounter{numberOfSections} % Counter for number of sections
\newtotcounter{numberOfQuestions} % Counter for number of questions

% EXAM STRUCTURE - SECTIONS AND QUESTIONS
\ifthenelse{\equal{\hasSections}{Yes} \OR \equal{\hasSections}{Y} \OR \equal{\hasSections}{yes} \OR \equal{\hasSections}{y} \OR \equal{\hasSections}{YES}}
% \hasSections THEN clause
{\newcommand{\sectionStart}{
    \setcounter{question}{0} % (Re)initiate questions counter
    \setcounter{figure}{0} % Initiate figure counter
    \setcounter{table}{0} % Initiate table counter
    \qformat{\ifthenelse{\equal{\thequestion}{1}}
        {\fbox{\parbox{16cm}{\textbf{\large SEZIONE \Alph{numberOfSections} \hfill\vrule depth 0.5em width 0pt \\ 
        Quesito \Alph{numberOfSections}\thequestion}} \hfill\vrule depth 1.5em width 0pt}}
        {\fbox{\parbox{16cm}{\large \textbf{Quesito \Alph{numberOfSections}\thequestion}} \hfill\vrule depth 0.8em width 0pt}}
    }
    \stepcounter{numberOfSections} % Increment sections counter
}}
% \hasSections ELSE clause
{   \setcounter{question}{0} % (Re)initiate questions counter
    \setcounter{figure}{0} % Initiate figure counter
    \setcounter{table}{0} % Initiate table counter
    \qformat{
        {\fbox{\parbox{16cm}{\large \textbf{Quesito \thequestion}} \hfill\vrule depth 0.8em width 0pt}}
    }
}

\newcommand{\callQuestion}[2]{ % Call new question
    \cleardoublepage \input{#1} % Start question in new page
    \addtocounter{pointsCounter}{#2} % Add points to counter
    \addtocounter{numberOfQuestions}{1} % Increment questions counter
    \stepcounter{figure} % Increment figure counter
    \stepcounter{table} % Increment table counter
}

\newcommand{\callFormulaSheet}[1]{
    \ifthenelse{\equal{\formulaSheet}{Yes} \OR \equal{\formulaSheet}{Y} \OR \equal{\formulaSheet}{yes} \OR \equal{\formulaSheet}{y} \OR \equal{\formulaSheet}{YES}}{\cleardoublepage \input{#1}}{}
}

\newcommand{\headingFormulaSheet}{
    \pagestyle{empty}
    \setlength{\textwidth}{180mm}
    \setlength{\textheight}{265mm}
    \setlength{\topmargin}{-20mm}
    \setlength\oddsidemargin{-12mm}
    \setlength\evensidemargin{-12mm}
    \setlength{\columnsep}{10mm}
    \setlength{\columnseprule}{0.4pt}
    \setlength{\jot}{5pt}
    \setlength\parindent{0pt}
    \twocolumn
    \fcolorbox{white}{myGrey}{\begin{minipage}{20em}
        \centering \vspace{4mm} {\Large \textbf{Formula Sheet}} \\ 
        \theCourse\ (\theCourseCode) \\
        {\small \theExamDate} \vspace{3mm} 
    \end{minipage}}
}

% TITLES AND FORMATS
\pointname{} % No title for marks
\pointsdroppedatright % Marks at right margin
\framedsolutions % Add frame to solutions
\addto\captionsenglish{ % Format Figure labels
    \renewcommand\figurename{Figure}}
\addto\captionsenglish{ % Format Table labels
    \renewcommand\tablename{Table}}
\makeatletter
\def\fnum@figure{\figurename~\Alph{numberOfSections}\thefigure}
\def\fnum@table{\tablename~\Alph{numberOfSections}\thetable}
\makeatother
\renewcommand{\solutiontitle}{
   \noindent\textbf{Soluzione:}\par\noindent} % Rename solutions

% COVER PAGE - DO NOT EDIT
\newcommand{\theStudentName}{Mario/a Rossi}
\newcommand{\theStudentMatriculation}{0542185}
\newcommand{\theStudentIDPath}{support/unipacard.png}

\usepackage[nomessages]{fp}% http://ctan.org/pkg/fp

\newcommand*{\ExtractFirstChar}[2]{%
    \StrRemoveBraces{#2}[\FirstChar]%
    \StrChar{\FirstChar}{#1}[\FirstChar]%
    \FirstChar}
\newcommand{\theStudentAlpha}{
    \ExtractFirstChar{4}{\theStudentMatriculation}}
\newcommand{\theStudentBeta}{
    \ExtractFirstChar{5}{\theStudentMatriculation}}
\newcommand{\theStudentGamma}{
    \ExtractFirstChar{6}{\theStudentMatriculation}}

%% Prepare Email
\usepackage{etoolbox}
\makeatletter
\newcommand\myemail[3]{%                %\newcommand\tpj@compose@mailto[3]{%
\edef\@tempa{mailto:#1?subject=#2 }%
\edef\@tempb{\expandafter\html@spaces\@tempa\@empty}%
\href{\@tempb}{#3}}
\catcode\%=11
\def\html@spaces#1 #2{#1%20\ifx#2\@empty\else\expandafter\html@spaces\fi#2}
\catcode\%=14
\makeatother
\newcommand{\theExamSubjectCode}{\theExamAcronym\theExamDate}
\newcommand{\theEmailSubject}{[\theExamSubjectCode] \theStudentName, \theStudentMatriculation}







%%%\usepackage[pdftex]{hyperref}

\usepackage{etoolbox}
\makeatletter
\let\HyField@calcorder\ltx@empty
\Hy@AtBeginDocument{%
  \if@filesw
    \immediate\write\@mainaux{%
      \string\providecommand\string\HyField@AuxAddToCalculationOrder[1]{}%
    }%
  \fi
  \def\HyField@AuxAddToCalculationOrder#1{%
    \xdef\HyField@calcorder{%
      \ifx\HyField@calcorder\@empty
      \else
        \HyField@calcorder
        \space
      \fi
      #1 0 R%
    }%
  }%
}
\def\HyField@@AddToFields#1{%
  \HyField@AfterAuxOpen{%
    \if@filesw
      \write\@mainaux{%
        \string\HyField@AuxAddToFields{#1}%
      }%
      \ifx\Fld@calculate@code\@empty
      \else
        \write\@mainaux{%
          \string\HyField@AuxAddToCalculationOrder{#1}%
        }%
      \fi
    \fi
  }%
}%
\patchcmd{\@Form}{/Fields[\HyField@afields]}{/CO[\HyField@calcorder]/Fields[\HyField@afields]}{}{}
\makeatother




\def\theExamDate{\today}



\newcommand{\afcover}
{
\MakeResultCmd{7}%\theafQuestionCnt}
%\meaning\afResultCmd % uncomment to inspect

\begin{coverpages}
\begin{flushright}
{\Large
\begin{Form}
\TextField[readonly=true, value=0, charsize=16pt, backgroundcolor=verylightgray, calculate={\afResultCmd}]{{\bf Risultato:~}}
\end{Form}
\\
}
{\em (a cura del docente)}
\end{flushright}
\hrule
%
\begin{flushleft}
\Large
\begin{tabular}{r p{11cm}}
Esame di: & {\bf \theCourse}\\
Tema: & {\bf \theThemeTitle}\\
Consegna: & {\bf \theExamDate}\\[20pt]
%
Gruppo: & {\bf \theGroupID}\\
Studenti: & {\bf \theStudentOne} (corrispondente)\\
\ifdefined\theStudentTwo & {\bf \theStudentTwo}\fi\\
\ifdefined\theStudentThree & {\bf \theStudentThree}\fi\\
\ifdefined\theStudentFour & {\bf \theStudentFour}\fi\\
\ifdefined\theStudentFive & {\bf \theStudentFive}\fi\\
\ifdefined\theStudentSix & {\bf \theStudentSix}\fi\\
\ifdefined\theStudentSeven & {\bf \theStudentSeven}\fi\\
\ifdefined\theStudentEight & {\bf \theStudentEight}\fi\\
\end{tabular}
\end{flushleft}
%
\vfill
%
% \begin{flushright}
% \tcbox[enhanced,drop shadow,colback=white]{\myemail{\theExamEmail}{\theEmailSubject}{{\em Prepara email: \theExamEmail} } } (se non automaticamente compilato, indicare come soggetto \\dell'email: ``\theEmailSubject''; \\ includere manualmente il PDF della soluzione)
% \end{flushright}
\vfill 
% \begin{flushleft}
% {\large \bf Dichiarazione di buona condotta durante lo svolgimento dell'esame}\\ \vspace{3mm} {\em Io sottoscritto(a) {\bf \theStudentName}, studente(essa) dell'Università degli Studi di Palermo, con numero di matricola {\bf\theStudentMatriculation}, dichiaro di aver svolto la soluzione qui sotto riportata in completa autonomia e senza l'uso di materiale non consentito.}
% \end{flushleft}
\begin{flushright}
\small Powered by {\bf MIRPALab.it}~\href{https://www.mirpalab.it}{\faHome}~\href{https://www.facebook.com/mirpalab/}{\faFacebookSquare}~\href{https://www.instagram.com/mirpalab/}{\faInstagram}
\end{flushright}
\end{coverpages}
\newpage
}

\ifthenelse{\equal{\withSolutions}{Yes} \OR \equal{\withSolutions}{Y} \OR \equal{\withSolutions}{yes} \OR \equal{\withSolutions}{y} \OR \equal{\withSolutions}{YES}}{\printanswers}{}

%
\def\afpt{0}

% AF: dynamic macros
%\newcommand\MakeArray[2]{\ifnumless{#2}{1}{\def#1{}}{\def#1{~0=0} \ifnumgreater{#2}{1}{\foreach \n in {1,...,#2} {\xdef#1{#1,~\n=\n}}}{}}}
\newcommand\MakeArray[2]{\ifnumless{#2}{1}{\def#1{}}{\def#1{0} \ifnumgreater{#2}{1}{\foreach \n in {1,...,#2} {\xdef#1{#1,\n}}}{}}}
%
\usepackage{pgffor}
\newcommand{\Repeat}[2]{\foreach \n in {1,...,#1}{#2}}
\newcommand\MakeResultCmd[1]{
\newcounter{ct}\setcounter{ct}{2}
\def\afResultCmd{event.value = this.getField('answerA').value} 
\Repeat{#1}{\xdef\afResultCmd{\afResultCmd+this.getField('answer\Alph{ct}').value} \addtocounter{ct}{1}} \xdef\afResultCmd{\afResultCmd;}}
%
\newcounter{afQuestionCnt}\setcounter{afQuestionCnt}{0}
\newcommand\afnotess[1]
{
\stepcounter{afQuestionCnt}
\begin{flushright}
\begin{Form}
\noindent
\MakeArray\MyArray#1
\TextField[name=tanswer\Alph{afQuestionCnt}, multiline=true, bordercolor=verylightgray, backgroundcolor=verylightgray, width=0.95\linewidth, height=1in, charsize=12pt, color=myRed]{{\bf Correzione} {\em (a cura del docente)}:\\[1pt]}\\
%\ChoiceMenu[radio, name=answer\Alph{afQuestionCnt}, value=0]{{\em Punteggio:~}}{\MyArray}
\TextField[name=answer\Alph{afQuestionCnt}, bordercolor=verylightgray, backgroundcolor=verylightgray, width=0.05\linewidth, height=0.05\linewidth, charsize=12pt, color=blue, align=1]{{\em Punteggio:~}}\\
\end{Form}
\end{flushright}
}
\def\afnotes{\afnotess{\afpt}}

% \def\psimbA{\afnotes\def\afpt{0}\droppoints}
% \def\psimbB{\afnotes\def\afpt{0}\droppoint}
% \def\psimbC{\afnotes\def\afpt{0}\droppoints}
% \def\psimbD{\afnotes\def\afpt{0}\droppoints}
% \def\psimbE{\afnotes\def\afpt{0}\droppoints}
% \def\psimbF{\afnotes\def\afpt{0}\droppoints}

% \def\pnumA{\afnotes\def\afpt{0}\droppoints}
% \def\pnumB{\afnotes\def\afpt{0}\droppoint}
% \def\pnumC{\afnotes\def\afpt{0}\droppoints}
% \def\pnumD{\afnotes\def\afpt{0}\droppoints}
% \def\pnumE{\afnotes\def\afpt{0}\droppoints}
% \def\pnumF{\afnotes\def\afpt{0}\droppoints}

\def\psimbA{}
\def\psimbB{}
\def\psimbC{}
\def\psimbD{}
\def\psimbE{}
\def\psimbF{}

\def\pnumA{}
\def\pnumB{}
\def\pnumC{}
\def\pnumD{}
\def\pnumE{}
\def\pnumF{}

\def\pteorA{}

% \def\finalComments
% {
% \begin{Form}
% \noindent
% \TextField[name=tCommenti, multiline=true, bordercolor=verylightgray, backgroundcolor=verylightgray, width=0.95\linewidth, height=1in, charsize=12pt, color=myRed]{{\bf Commenti generali} {\em (a cura del docente)}:\\[1pt]}\\
% \end{Form}
% }

\def\finalComments{}




\def\afTheme{%
%
\section*{Tema {\em \theThemeTitle}}
\begin{multicols}{2} \noindent\theThemeDescription \thePicture \end{multicols}
\vfill\theModel\vfill 
\begin{multicols}{2} \theVariables\theParameters \end{multicols}
\theRequirements
\vfill \newpage
%
}



\def\afSubmission{%

\section*{Istruzioni per la consegna}
{\small 

\noindent Ogni tema progettuale si propone di far sperimentare allo studente, sia pure in modo semplificato, i diversi momenti di analisi e progettazione, sottesi alla soluzione di un problema di controllo; ciò avviene, per ovvie ragioni, attraverso uno studio teorico e solo simulato, come è tipico nei corsi di controlli automatici di base, ma rappresenta comunque un momento fondamentale che precede qualsiasi realizzazione completa e sperimentale. Le Tabelle~$1$ e~$2$ indicano il significato delle variabili e dei parametri del modello~(\ref{eq:model}), mentre le prima due righe della Tabella~$3$ specificano le ipotesi operative di equilibrio, mentre le restanti righe i requisiti di progetto.\\

La consegna del progetto avviene inviando un'email all'indirizzo di posta elettronica~\textit{\theExamEmail}, con oggetto ``[\theGroupID] Consegna'' e con allegato il PDF generato da questo documento. La consegna può essere fatta solo dal corrispondente del gruppo.\\

Dal momento della consegna non sarà più possibile modificare il seguente progetto overleaf, il quale dovrà includere almeno i seguenti contenuti:
\begin{itemize}
%
\item i file sorgenti \LaTeX{} del PDF (main.tex, theme.tex, 1*.tex, 2*tex, 3*.tex);
%
\item una cartella \textit{img} contenente le immagini e tutti i grafici inseriti nel testo, che devono essere riportati in formato vettoriale e contenere annotazioni adeguate e leggibili;
\item una cartella \textit{support} contenente (1) gli script (estensione .m) necessari per inizializzare e lanciare la simulazione, (2) i salvataggi delle sessioni di lavoro in \emph{sisotool}, e (3) i due diagrammi \emph{Simulink} (estensione .slx) dei sistemi approssimato lineare e non lineare, controllati (eventualmente con il doppio anello di retroazione), correttamente funzionanti e tali da riprodurre i risultati riportati nel documento descrittivo. 
\end{itemize}
\vfill \newpage

}

}




\def\theGroupID{CA22-21}
\def\theStudentOne{Bonura Davide, 0713790}
\def\theStudentTwo{La Rosa Mazza Nicolò, 0712006}
\def\theStudentThree{Bova Gabriele, 0711771}
\def\theStudentFour{Cimino Giuseppe, 0713924}
\def\theStudentFive{Carini Riccardo, 0713323}
\def\theStudentSix{Aldegheri Amedeo Dante, 0711827}

\def\theThemeFile{theme.tex}
%\def\theThemeFile{theme-lanciatore.tex}
%\def\theThemeFile{theme-catamarano.tex}
\def\theThemeFile{theme-sottomarino.tex}
%\def\theThemeFile{theme-biomassa.tex}
%\def\theThemeFile{theme-predapredatori.tex}



\input{\theThemeFile}

%
%
%

\begin{document}

\afcover
\afTheme
\afSubmission

%
%
%

\section{Parte Simbolica}

Si risponda ai seguenti quesiti in modo simbolico:\\

\begin{parts}	

\part[\afpt] Usando le condizioni (H1) nella Tabella~$3$, trovare le configurazioni di equilibrio del sistema.
\begin{solution} 
Al fine di trovare il punto di equilibrio del sistema, è necessario considerare nulle tutte le derivate del sistema.\\
Poniamo quindi le derivate dell'equazione \eqref{eq:model} pari a zero e otteniamo il seguente risultato:
\begin{displaymath}
\left\{
\begin{array}{l}
0=\bar F \cos (\bar\vartheta)-\bar F_L \operatorname{sin}(\bar\vartheta)-\bar B \\
0=0-k_{\vartheta} \operatorname{sin}(\bar\vartheta)-\bar{F} b \\
y=\bar{z}-b \operatorname{sin}(\bar\vartheta)
\end{array}
\right.
\end{displaymath}
Svolgendo i calcoli otteniamo:
\begin{displaymath}
\left\{
\begin{array}{l}
\bar F_L= \frac{\bar F \cos(\bar\vartheta)-\bar B}{\sin(\bar\vartheta)} \\
\bar F = -\frac{k_\vartheta \sin(\bar\vartheta)}{b} \\
\bar z = y + b\sin(\bar\vartheta)
\end{array}
\right.
\end{displaymath}
Sostituendo $\bar F$ in $\bar F_L$:
\begin{displaymath}
\left\{
\begin{array}{l}
\bar F_L= \frac{(-\frac{k_\vartheta \sin(\bar\vartheta)}{b}) \cos(\bar\vartheta)-\bar B}{\sin(\bar\vartheta)} \\
\bar F = -\frac{k_\vartheta \sin(\bar\vartheta)}{b} \\
\bar z = y + b\sin(\bar\vartheta) \\
\end{array}
\right.
\end{displaymath}
Sostituendo il valore $\bar\vartheta = \frac{\pi}{6}$:
\begin{displaymath}
\left\{
\begin{array}{l}
\bar F_L= 2(-\frac{\sqrt3}{4}\frac{k_\vartheta}{b}-\bar B)\\
\bar F = -\frac{k_\vartheta}{2b} \\
\bar z = y + \frac{b}{2} \\
\end{array}
\right.
\end{displaymath}
\end{solution}
\psimbA

\part[\afpt] Assumendo piccoli spostamenti dello stato del sistema dagli equilibri trovati sopra, determinare il modello linearizzato per le variabili traslate.
\begin{solution} 
Vengono definiti l'ingresso F e l’uscita del sistema $y=z-bsen(\theta)$. \\Inoltre, viene definito il vettore delle variabili di stato $x=[x_1,x_2,x_3,x_4]$:
\begin{displaymath}
\left\{
\begin{array}{l}
x_{1}\triangleq z \\ 
x_{2}\triangleq \vartheta \\ 
x_{3}\triangleq z' \\
x_{4}\triangleq \vartheta'\\
\end{array}
\right.
\end{displaymath}
Definisco il vettore degli ingressi di controllo $u=F$ \\
Il sistema assume la seguente forma:
\begin{displaymath}
\left\{
\begin{array}{l}
x_{1}'\triangleq z'=x_3 \\
x_{2}'\triangleq \vartheta' = x_4 \\
x_{3}'\triangleq z'' \\
x_{4}'\triangleq\vartheta''
\end{array}
\right.
\end{displaymath}
Per ricavare z'' e $\vartheta$'' è possibile portare in forma normale le equazioni del testo:
\begin{displaymath}
\begin{array}{l}
z''=\frac{Fcos(\vartheta)}{M}-\frac{k_{z}z'cos(\vartheta)}{M}-\frac{F_{L}sin(\vartheta)}{M}-\frac{B}{M}=\frac{ucos(x_2)}{M}-\frac{k_{z}x_3cos(x_2)}{M}-\frac{F_Lsin(x_2)}{M}-\frac{B}{M} \\
\vartheta''=-\frac{k_{r}\vartheta'}{J}-\frac{k_{\vartheta}sin(\vartheta)}{J}-\frac{Fb}{J}=-\frac{k_{r}x_{4}}{J}-\frac{k_{\vartheta}sin(x_2)}{J}-\frac{ub}{J}
\end{array}
\end{displaymath}
Notiamo che il sistema non è lineare e procediamo a linearlizzarlo tramite linearizzazione approssimata. \\
Siano gli equilibri:
\begin{displaymath}
\begin{array}{c}
\bar x = [\bar z,\bar\vartheta,0,0] = [\bar x_1,\bar x_2,\bar x_3,\bar x_4] \\
\bar u = \bar F
\end{array}
\end{displaymath}
\begin{comment}
Avremo quindi:
\begin{displaymath}
\begin{array}{l}
\delta x_1=x_1-\bar{x_1} \\
\delta x_2=x_2 - \bar x_2 \\
%\delta x_3=x_3 \\
%\delta x_4=x_4 \\
\delta u=u-\bar u \\
\end{array}
\end{displaymath}
\end{comment}
Sia $Dx = f(x,u)$. Sviluppando in serie di Taylor e troncando al primo ordine, si ha
\begin{displaymath}
\begin{array}{l}
f(x,u) \approx f(\bar x,\bar u) + \frac{\partial f}{\partial x}\mid_{\binom{x=\bar{x}}{u=\bar{u}}}(x-\bar x) + \frac{\partial f}{\partial u}\mid_{\binom{x=\bar{x}}{u=\bar{u}}}(u-\bar u)
\end{array}
\end{displaymath}
Che quindi si semplifica in:
\begin{displaymath}
\begin{array}{l}
f(x,u) \approx \frac{\partial f}{\partial x}\mid_{\binom{x=\bar{x}}{u=\bar{u}}}\delta x + \frac{\partial f}{\partial u}\mid_{\binom{x=\bar{x}}{u=\bar{u}}}\delta u
\end{array}
\end{displaymath}
Definendo:
\begin{displaymath}
\begin{array}{l}
A=\frac{\partial f}{\partial x}\mid_{\binom{x=\bar{x}}{u=\bar{u}}} \\
B=\frac{\partial f}{\partial u}\mid_{\binom{x=\bar{x}}{u=\bar{u}}} \\
\end{array}
\end{displaymath}
Si osserva che A e B sono matrici rispettivamente n×n e n×m a valori costanti. Introducendo nuove variabili di stato $x-\bar x$, e nuovi ingressi $u-\bar u$ come differenza rispetto ai valori nominali, si osserva facilmente che alle condizioni di equilibrio in tempo continuo $(Dx=\bar x'=f(\bar x;\bar u)=0)$ il sistema risulta una approssimazione al primo ordine del sistema dato. \\
Otteniamo il sistema lineare $D\bar x = A \bar x + B\bar u$ che in forma matriciale sarà: 
\begin{displaymath}
\begin{array}{l}
\left(\begin{array}{c}
x_{1}'\\
x_{2}'\\
x_{3}'\\
x_{4}'
\end{array}\right)=\left(\begin{array}{cccc}
\frac{\partial x_{1}'}{dx_{1}} & \frac{\partial x_{1}'}{dx_{2}} & \frac{\partial x_{1}'}{dx_{3}} & \frac{\partial x_{1}'}{dx_{4}}\\
\frac{\partial x_{2}'}{dx_{1}} & \frac{\partial x_{2}'}{dx_{2}} & \frac{\partial x_{2}'}{dx_{3}} & \frac{\partial x_{2}'}{dx_{4}}\\
\frac{\partial x_{3}'}{dx_{1}} & \frac{\partial x_{3}'}{dx_{2}} & \frac{\partial x_{3}'}{dx_{3}} & \frac{\partial x_{3}'}{dx_{4}}\\
\frac{\partial x_{4}'}{dx_{1}} & \frac{\partial x_{4}'}{dx_{2}} & \frac{\partial x_{4}'}{dx_{3}} & \frac{\partial x_{4}'}{dx_{4}}
\end{array}\right)_{\binom{x=\bar{x}}{u=\bar{u}}}\left(\begin{array}{c}
x_{1}\\
x_{2}\\
x_{3}\\
x_{4}
\end{array}\right)+\left(\begin{array}{cc}
\frac{\partial x_{1}'}{du} \\
\frac{\partial x_{2}'}{du} \\
\frac{\partial x_{3}'}{du} \\
\frac{\partial x_{4}'}{du} \\
\end{array}\right)_{\binom{x=\bar{x}}{u=\bar{u}}}\left(\begin{array}{c}
u
\end{array}\right)

\end{array}
\end{displaymath}
Svolgendo i calcoli otteniamo:
\begin{displaymath}
\left(\begin{array}{c}
x_{1}'\\
x_{2}'\\
x_{3}'\\
x_{4}'
\end{array}\right)=\left(\begin{array}{cccc}
0 & 0 & 1 & 0\\
0 & 0 & 0 & 1\\
0 & -\frac{\bar usin(\bar{x_{2}})}{M}+\frac{k_{z}\bar{x_{3}}sin(\bar{x_{2}})}{M}-\frac{\bar{F_L}cos(\bar{x_{2}})}{M} & -\frac{k_{z}cos(\bar{x_{2}})}{M} & 0\\
0 & -\frac{k_{\vartheta}cos(\bar{x_{2}})}{J} & 0 & -\frac{k_{r}}{J}
\end{array}\right)\left(\begin{array}{c}
x_{1}\\
x_{2}\\
x_{3}\\
x_{4}
\end{array}\right)+ 
\end{displaymath}
\begin{displaymath}
+\left(\begin{array}{c}
0\\
0\\
\frac{cos(\bar{x_{2}})}{M}\\
-\frac{b}{J}
\end{array}\right)\left(\begin{array}{c}
u\end{array}\right)
\end{displaymath}
Osserviamo che per come abbiamo definito il vettore di stato, la matrice A presenta una forma a blocchi: in alto a sinistra c'è una sottomatrice 2x2 nulla, in alto a destra c'è una sottomatrice 2x2 identità, in basso a sinistra una sottomatrice 2x2 triangolare superiore e in basso a destra una sottomatrice 2x2 diagonale. Questa struttura potrebbe essere utile per compattare le informazioni (ad esempio richiedendo meno banda per trasmetterle). \\
Da qui in poi per comodità poniamo:
\begin{displaymath}
\begin{array}{l}
\alpha=-\frac{k_{z}cos(\bar{x_{2}})}{M}<0 \\
\beta=-\frac{\bar usin(\bar{x_{2}})} {M}+\frac{k_{z}\bar{x_{2}}sin(\bar{x_{2}})}{M}-\frac{\bar{F_L}cos(\bar{x_{2}})}{M}>0 \\
\gamma=-\frac{k_{\vartheta}cos(\bar{x_{2}})}{J}<0 \\
\varepsilon=-\frac{k_{r}}{J}<0 \\
\end{array}
\end{displaymath}
ottenendo: 
\begin{displaymath}
\left(\begin{array}{c}
x_{1}'\\
x_{2}'\\
x_{3}'\\
x_{4}'
\end{array}\right)=\left(\begin{array}{cccc}
0 & 0 & 1 & 0\\
0 & 0 & 0 & 1\\
0 & \beta & \alpha & 0\\
0 & \gamma & 0 & \varepsilon
\end{array}\right)\left(\begin{array}{c}
x_{1}\\
x_{2}\\
x_{3}\\
x_{4}
\end{array}\right)+\left(\begin{array}{c}
0\\
0\\
\frac{cos(\bar{x_{2}})}{M}\\
-\frac{b}{J}
\end{array}\right)\left(\begin{array}{c}
u\end{array}\right)
\end{displaymath}
Se vi è una mappa di uscita $y=h(x,u)$, si potrà linearizzare anch’essa in modo simile, ponendo $y =y(x,u) - y(\bar x, \bar u)=Cx+Du$, dove le matrici C e D sono esprimibili come:
\begin{displaymath}
\begin{array}{l}
C=\frac{\partial y}{\partial x}\mid_{\binom{x=\bar{x}}{u=\bar{u}}} \\
D=\frac{\partial y}{\partial u}\mid_{\binom{x=\bar{x}}{u=\bar{u}}} \\
\end{array}
\end{displaymath}
In particolare l'espressione dell'uscita viene scritta come:
\begin{displaymath}
y=\left(\begin{array}{cccc}
\frac{\partial y}{dx_{1}} & \frac{\partial y}{dx_{2}} & \frac{\partial y}{dx_{3}} & \frac{\partial y}{dx_{4}}\end{array}\right)_{\binom{x=\bar{x}}{u=\bar{u}}}\left(\begin{array}{c}
x_{1}\\
x_{2}\\
x_{3}\\
x_{4}
\end{array}\right)+\left(\begin{array}{cc}
\frac{\partial y}{du}\end{array}\right)_{\binom{x=\bar{x}}{u=\bar{u}}}\left(\begin{array}{c}
u
\end{array}\right)
\end{displaymath}
Svolgendo i calcoli otteniamo:
\begin{displaymath}
y=\left(\begin{array}{cccc}
1 & -bcos(\bar{x_{2}}) & 0 & 0\end{array}\right)_{\binom{x=\bar{x}}{u=\bar{u}}}\left(\begin{array}{c}
x_{1}\\
x_{2}\\
x_{3}\\
x_{4}
\end{array}\right)+\left(\begin{array}{cc}
0 \end{array}\right)_{\binom{x=\bar{x}}{u=\bar{u}}}\left(\begin{array}{c}
u
\end{array}\right)
\end{displaymath}


\end{solution}
\psimbB

\part[\afpt] Indicando con $x_0$ il generico stato iniziale, determinare l'evoluzione libera del modello ed esprimerla come $M(A,x_0)\, m(t)$, con $m(t)$ il vettore dei modi.
\begin{solution} 
Calcoliamo il polinomio caratteristico:
\begin{displaymath}
\pi_{A}(\lambda)=det(A-\lambda I)=det\left(\begin{array}{cccc}
-\lambda & 0 & 1 & 0\\
0 & -\lambda & 0 & 1\\
0 & \beta & \alpha-\lambda & 1\\
0 & \gamma & 0 & \varepsilon-\lambda
\end{array}\right)=
\end{displaymath}
\begin{scriptsize}
\begin{displaymath}
=(-1)^{1+1}(-\lambda)det\left(\begin{array}{ccc}
-\lambda & 0 & 1\\
\beta & \alpha-\lambda & 1\\
\gamma & 0 & \varepsilon-\lambda
\end{array}\right)=(-1)^{1+1}(-\lambda)(-1)^{2+2}(\alpha-\lambda)det\left(\begin{array}{cc}
-\lambda & 1\\
\gamma & \varepsilon-\lambda
\end{array}\right)= 
\end{displaymath}
\end{scriptsize}
\begin{displaymath}
=(-\lambda)(\alpha-\lambda)(-\lambda(\varepsilon-\lambda)-\gamma)=(-\lambda)(\alpha-\lambda)(\lambda^{2}-\varepsilon\lambda-\gamma)
\end{displaymath}
Scriviamo l'equazione caratteristica:
\begin{displaymath}
(-\lambda)(\alpha-\lambda)(\lambda^{2}-\varepsilon\lambda-\gamma) = 0
\end{displaymath}
Otteniamo i seguenti autovalori reali:
\begin{displaymath}
\begin{array}{cc}
\lambda_1=0 \\
\lambda_2=\alpha<0
\end{array}
\end{displaymath}
Per risolvere l'equazione $(\lambda^{2}-\varepsilon\lambda-\gamma)$ calcoliamo $\Delta=\varepsilon^2+4\gamma$ e notiamo che viene negativo.
Applicando la regola di Cartesio, poichè $\varepsilon$ e $\gamma$ sono entrambi negativi, l'equazione presenta due permanenze che corrispondono a due soluzioni complesse coniugate con parte reale negativa. \\
Per calcolare l'evoluzione libera utilizziamo $x(t)=e^{At}x_{0}$. \\
Dopo aver verificato che A sia diagonalizzabile, procediamo nel seguente modo:
\begin{displaymath}
e^{At}=Te^{\Lambda t}T^{-1}
\end{displaymath}
Sia $T=(v_1|v_2|v_3|v_4)$, dove $v_1,v_2,v_3,v_4$ sono gli autovettori. \\
Calcolando gli autovettori e mettendoli per colonna otteniamo la seguente matrice:
\begin{displaymath}
T=\left(\begin{array}{cccc}
1 & \frac{1}{\alpha} & -\frac{2\beta(\sqrt{\Delta}+\varepsilon)}{\gamma(\sqrt{\Delta}-\varepsilon)(2\alpha+\sqrt{\Delta}-\varepsilon)} & \frac{2\beta(\sqrt{\Delta}-\varepsilon)}{\gamma(\sqrt{\Delta}+\varepsilon)(-2\alpha+\sqrt{\Delta}+\varepsilon)}\\
0 & 0 & -\frac{\sqrt{\Delta}+\varepsilon}{2\gamma} & \frac{\sqrt{\Delta}-\varepsilon}{2\gamma} \\
0 & 1 & \frac{\beta(\sqrt{\Delta}+\varepsilon)}{\gamma(2\alpha+\sqrt{\Delta}-\varepsilon)}  & \frac{\beta(\sqrt{\Delta}-\varepsilon)}{\gamma(-2\alpha+\sqrt{\Delta}+\varepsilon)}\\
0 & 0 & 1 & 1\\
\end{array}\right)
\end{displaymath}
Invertendo la matrice appena calcolata otteniamo:
\begin{displaymath}
T^{-1}=\left(\begin{array}{cccc}
1 & -\frac{\beta\varepsilon}{\alpha\gamma} & -\frac{1}{\alpha} & \frac{\beta}{\alpha\gamma}\\
0 & \frac{\beta(\varepsilon-\alpha)}{-\alpha^{2}+\alpha\varepsilon+\gamma} & 1 & \frac{-\beta}{-\alpha^{2}+\alpha\varepsilon+\gamma}\\
0 & -\frac{\gamma}{\sqrt{\Delta}} & 0 & \frac{1}{2}-\frac{\varepsilon}{2\sqrt{\Delta}}\\
0 & \frac{\gamma}{\sqrt{\Delta}} & 0 & \frac{1}{2}(\frac{\varepsilon}{\sqrt{\Delta}}+1)
\end{array}\right)
\end{displaymath}
Per ricavare $\Lambda$ è sufficiente scrivere una matrice diagonale posizionando sulla diagonale principale gli autovalori $\lambda_i$:
\begin{displaymath}
\Lambda=\left(\begin{array}{cccc}
0 & 0 & 0 & 0\\
0 & \alpha & 0 & 0\\
0 & 0 & \frac{1}{2}(\varepsilon-\sqrt{\Delta}) & 0\\
0 & 0 & 0 & \frac{1}{2}(\varepsilon+\sqrt{\Delta})
\end{array}\right)
\end{displaymath}
E poiché $\Lambda$ è diagonale per costruzione possiamo subito ricavare:
\begin{displaymath}
e^{\Lambda t}=\left(\begin{array}{cccc}
e^{\lambda_{1}t} & 0 & 0 & 0\\
0 & e^{\lambda_{2}t} & 0 & 0\\
0 & 0 & e^{\lambda_{3}t} & 0\\
0 & 0 & 0 & e^{\lambda_{4}t}
\end{array}\right)
\end{displaymath}
A questo punto, come detto prima, vale la relazione:
\begin{displaymath}
x(t)=e^{At}x_{0}=Te^{\Lambda t}T^{-1}x_{0}
\end{displaymath}
E svolgendo i calcoli otteniamo:
\begin{displaymath}
Te^{\Lambda t}T^{-1}=\left(\begin{array}{cccc}
a_{11} & a_{12} & a_{13} & a_{14}\\
a_{21} & a_{22} & a_{23} & a_{24}\\
a_{31} & a_{32} & a_{33} & a_{34}\\
a_{41} & a_{42} & a_{43} & a_{44}\\
\end{array}\right)
\end{displaymath}
dove:
\begin{tiny}
\begin{displaymath}
\begin{array}{c}
a_{11}=e^{\lambda_{1}t} \\ \\
a_{12}=e^{\lambda_{1}t}(-\frac{\beta\varepsilon}{\alpha\gamma})+\frac{1}{\alpha}e^{\lambda_{2}t}(\frac{\beta(\varepsilon-\alpha)}{-\alpha^{2}+\alpha\varepsilon+\gamma})+\frac{2\beta(\sqrt{\Delta}+\varepsilon)}{\gamma(\sqrt{\Delta}-\varepsilon)(2\alpha+\sqrt{\Delta}-\varepsilon)}e^{\lambda_{3}t}(\frac{\gamma}{\sqrt{\Delta}})+\frac{2\beta(\sqrt{\Delta}-\varepsilon)}{\gamma(\sqrt{\Delta}+\varepsilon)(-2\alpha+\sqrt{\Delta}+\varepsilon)}e^{\lambda_{4}t}(\frac{\gamma}{\sqrt{\Delta}}) \\ \\
a_{13}=\frac{1}{\alpha}(e^{\lambda_{2}t}-e^{\lambda_{1}t}) \\ \\
a_{14}=e^{\lambda_{1}t}\frac{\beta}{\alpha\gamma}+\frac{1}{\alpha}e^{\lambda_{2}t}(\frac{-\beta}{-\alpha^{2}+\alpha\varepsilon+\gamma})-\frac{2\beta(\sqrt{\Delta}+\varepsilon)}{\gamma(\sqrt{\Delta}-\varepsilon)(2\alpha+\sqrt{\Delta}-\varepsilon)}e^{\lambda_{3}t}(\frac{1}{2}-\frac{\varepsilon}{2\sqrt{\Delta}})+\frac{2\beta(\sqrt{\Delta}-\varepsilon)}{\gamma(\sqrt{\Delta}+\varepsilon)(-2\alpha+\sqrt{\Delta}+\varepsilon)}e^{\lambda_{4}t}(\frac{1}{2}(\frac{\varepsilon}{\sqrt{\Delta}}+1)) \\ \\
a_{21}=0 \\ \\
a_{22}=\frac{\sqrt{\Delta}+\varepsilon}{2\gamma}e^{\lambda_{3}t}(\frac{\gamma}{\sqrt{\Delta}})+\frac{\sqrt{\Delta}-\varepsilon}{2\gamma}e^{\lambda_{4}t}(\frac{\gamma}{\sqrt{\Delta}}) \\ \\
a_{23}=0 \\ \\
a_{24}=-\frac{\sqrt{\Delta}+\varepsilon}{2\gamma}e^{\lambda_{3}t}(\frac{1}{2}-\frac{\varepsilon}{2\sqrt{\Delta}})+\frac{\sqrt{\Delta}-\varepsilon}{2\gamma}e^{\lambda_{4}t}(\frac{1}{2}(\frac{\varepsilon}{\sqrt{\Delta}}+1)) \\ \\
a_{31}=0 \\ \\
a_{32}=e^{\lambda_{2}t}(\frac{\beta(\varepsilon-\alpha)}{-\alpha^{2}+\alpha\varepsilon+\gamma})-\frac{\beta(\sqrt{\Delta}+\varepsilon)}{\gamma(2\alpha+\sqrt{\Delta}-\varepsilon)}e^{\lambda_{3}t}(\frac{\gamma}{\sqrt{\Delta}})+\frac{\beta(\sqrt{\Delta}-\varepsilon)}{\gamma(-2\alpha+\sqrt{\Delta}+\varepsilon)}e^{\lambda_{4}t}(\frac{\gamma}{\sqrt{\Delta}}) \\ \\
a_{33}=e^{\lambda_{2}t} \\ \\
a_{34}=e^{\lambda_{2}t}(\frac{-\beta}{-\alpha^{2}+\alpha\varepsilon+\gamma})+\frac{\beta(\sqrt{\Delta}+\varepsilon)}{\gamma(2\alpha+\sqrt{\Delta}-\varepsilon)}e^{\lambda_{3}t}(\frac{1}{2}-\frac{\varepsilon}{2\sqrt{\Delta}})+\frac{\beta(\sqrt{\Delta}-\varepsilon)}{\gamma(-2\alpha+\sqrt{\Delta}+\varepsilon)}e^{\lambda_{4}t}(\frac{1}{2}(\frac{\varepsilon}{\sqrt{\Delta}}+1)) \\ \\
a_{41}=0 \\ \\
a_{42}=e^{\lambda_{3}t}(-\frac{\gamma}{\sqrt{\Delta}})+e^{\lambda_{4}t}(\frac{\gamma}{\sqrt{\Delta}}) \\ \\
a_{43}=0 \\ \\
a_{44}=e^{\lambda_{3}t}(\frac{1}{2}-\frac{\varepsilon}{2\sqrt{\Delta}})+e^{\lambda_{4}t}(\frac{1}{2}(\frac{\varepsilon}{\sqrt{\Delta}}+1)) \\ \\
\end{array}
\end{displaymath}
\end{tiny}
Infine moltiplicando per $x_{0}$ otteniamo:
\begin{displaymath}
x(t)=Te^{\Lambda t}T^{-1}x_{0}=\left(\begin{array}{c}
b_{1}\\
b_{2}\\
b_{3}\\
b_{4}
\end{array}\right)
\end{displaymath}
dove:
\begin{tiny}
\begin{displaymath}
\begin{array}{l}
b_{1}={\color{red}e^{\lambda_{1}t}x_{0_{1}}+e^{\lambda_{1}t}(-\frac{\beta\varepsilon}{\alpha\gamma})x_{0_{2}}}{\color{green}+\frac{1}{\alpha}e^{\lambda_{2}t}(\frac{\beta(\varepsilon-\alpha)}{-\alpha^{2}+\alpha\varepsilon+\gamma})x_{0_{2}}}{\color{blue}+\frac{2\beta(\sqrt{\Delta}+\varepsilon)}{\gamma(\sqrt{\Delta}-\varepsilon)(2\alpha+\sqrt{\Delta}-\varepsilon)}e^{\lambda_{3}t}(\frac{\gamma}{\sqrt{\Delta}})x_{0_{2}}} \color{black}+\\ \\ {\color{purple}+ \frac{2\beta(\sqrt{\Delta}-\varepsilon)}{\gamma(\sqrt{\Delta}+\varepsilon)(-2\alpha+\sqrt{\Delta}+\varepsilon)}e^{\lambda_{4}t}(\frac{\gamma}{\sqrt{\Delta}})x_{0_{2}}+} 
{\color{green}+\frac{1}{\alpha}e^{\lambda_{2}t}x_{0_{3}}}{\color{red}-\frac{1}{\alpha}e^{\lambda_{1}t}x_{0_{3}}}{\color{red}+e^{\lambda_{1}t}\frac{\beta}{\alpha\gamma}x_{0_{4}}}{\color{green}+\frac{1}{\alpha}e^{\lambda_{2}t}(\frac{-\beta}{- \alpha^{2}+\alpha\varepsilon+\gamma})x_{0_{4}}} + \\ \\ {\color{blue}-\frac{2\beta(\sqrt{\Delta}+\varepsilon)}{\gamma(\sqrt{\Delta}-\varepsilon)(2\alpha+\sqrt{\Delta}-\varepsilon)}e^{\lambda_{3}t}(\frac{1}{2}-\frac{\varepsilon}{2\sqrt{\Delta}})x_{0_{4}}}{\color{purple}+\frac{2\beta(\sqrt{\Delta}-\varepsilon)}{\gamma(\sqrt{\Delta}+\varepsilon)(-2\alpha+\sqrt{\Delta}+\varepsilon)}e^{\lambda_{4}t}(\frac{1}{2}(\frac{\varepsilon}{\sqrt{\Delta}}+1))x_{0_{4}}} \\ \\
b_{2}={\color{blue}\frac{\sqrt{\Delta}+\varepsilon}{2\gamma}e^{\lambda_{3}t}(\frac{\gamma}{\sqrt{\Delta}})x_{0_{2}}}{\color{purple}+\frac{\sqrt{\Delta}-\varepsilon}{2\gamma}e^{\lambda_{4}t}(\frac{\gamma}{\sqrt{\Delta}})x_{0_{2}}}{\color{blue}-\frac{\sqrt{\Delta}+\varepsilon}{2\gamma}e^{\lambda_{3}t}(\frac{1}{2}-\frac{\varepsilon}{2\sqrt{\Delta}})x_{0_{4}}}{\color{purple}+\frac{\sqrt{\Delta}-\varepsilon}{2\gamma}e^{\lambda_{4}t}(\frac{1}{2}(\frac{\varepsilon}{\sqrt{\Delta}}+1))x_{0_{4}}} \\ \\
\end{array}
\end{displaymath}
\begin{displaymath}
\begin{array}{l}
b_{3}={\color{green}e^{\lambda_{2}t}(\frac{\beta(\varepsilon-\alpha)}{-\alpha^{2}+\alpha\varepsilon+\gamma})x_{0_{2}}}{\color{blue}-\frac{\beta(\sqrt{\Delta}+\varepsilon)}{\gamma(2\alpha+\sqrt{\Delta}-\varepsilon)}e^{\lambda_{3}t}(\frac{\gamma}{\sqrt{\Delta}})x_{0_{2}}}{\color{purple}+\frac{\beta(\sqrt{\Delta}-\varepsilon)}{\gamma(-2\alpha+\sqrt{\Delta}+\varepsilon)}e^{\lambda_{4}t}(\frac{\gamma}{\sqrt{\Delta}})x_{0_{2}}}{\color{green}+e^{\lambda_{2}t}x_{0_{3}}}{\color{green}+} \\ \\
{\color{green}+}{\color{green}e^{\lambda_{2}t}(\frac{-\beta}{-\alpha^{2}+\alpha\varepsilon+\gamma})x_{0_{4}}}{\color{blue}+\frac{\beta(\sqrt{\Delta}+\varepsilon)}{\gamma(2\alpha+\sqrt{\Delta}-\varepsilon)}e^{\lambda_{3}t}(\frac{1}{2}-\frac{\varepsilon}{2\sqrt{\Delta}})x_{0_{4}}}{\color{purple}+\frac{\beta(\sqrt{\Delta}-\varepsilon)}{\gamma(-2\alpha+\sqrt{\Delta}+\varepsilon)}e^{\lambda_{4}t}(\frac{1}{2}(\frac{\varepsilon}{\sqrt{\Delta}}+1))x_{0_{4}}} \\ \\
b_{4}={\color{blue}e^{\lambda_{3}t}(-\frac{\gamma}{\sqrt{\Delta}})x_{0_{2}}}{\color{purple}+e^{\lambda_{4}t}(\frac{\gamma}{\sqrt{\Delta}})x_{0_{2}}}{\color{blue}+e^{\lambda_{3}t}(\frac{1}{2}-\frac{\varepsilon}{2\sqrt{\Delta}})x_{0_{4}}}{\color{purple}+e^{\lambda_{4}t}(\frac{1}{2}(\frac{\varepsilon}{\sqrt{\Delta}}+1))x_{0_{4}}} \\
\end{array}
\end{displaymath}
\end{tiny}
Pur essendo presenti nella soluzione $x(t)$ quantità complesse come autovalori $\lambda_{i}$ complessi e $\sqrt{\Delta}$ con $\Delta<0$, sappiamo dalla teoria che autovalori e autovettori complessi appaiono sempre insieme al loro coniugato e poiché A è a valori reali (come in tutti i casi affrontati durante il corso) questo fa sì che $x(t)$ sia comunque complessivamente a valori reali. 
Abbiamo verificato in forma numerica con Matlab la correttezza di tale risultato.
\\\\In alternativa si poteva applicare l'altro metodo visto a lezione che prevede l'utilizzo di un cambiamento di coordinate al fine di trasformare per similitudine la matrice diagonalizzabile (sui complessi) A in una matrice reale diagonale a blocchi. In questo modo si poteva arrivare direttamente al risultato in forma reale. \\
\\ Dalla teoria dell'analisi modale (ricerca e studio delle funzioni che compaiono nella soluzione libera) sappiamo che: \\
1) Autovalori $\lambda$ reali generano modi esponenziali semplici del tipo $e^{\lambda t}$ \\
2) Autovalori $\lambda$ complessi generano modi oscillanti del tipo $e^{\sigma t}\cos(\omega t)$ e $e^{\sigma t}\sin(\omega t)$\\
Questo ci permette di scrivere il vettore dei modi, indispensabile per descrivere il modo con cui il modello evolve:
\begin{displaymath}
m(t)=\left(\begin{array}{c}
e^{\lambda_{1}t}\\
e^{\lambda_{2}t}\\
e^{\sigma t}\cos(\omega t)\\
e^{\sigma t}\sin(\omega t)
\end{array}\right)=\left(\begin{array}{c}
1 \\
e^{\alpha t} \\
e^{\sigma t}\cos(\omega t)\\
e^{\sigma t}\sin(\omega t)
\end{array}\right)
=\left(\begin{array}{c}
1\\
e^{\alpha t} \\
e^{\frac{\varepsilon}{2}t}\cos(\operatorname{Im}(\frac{\sqrt{\Delta}}{2}))\\
e^{\frac{\varepsilon}{2}t}\sin(\operatorname{Im}(\frac{\sqrt{\Delta}}{2}))
\end{array}\right)
\end{displaymath}
A questo punto dall'analisi matematica ricordiamo che: 
\begin{displaymath}
\begin{array}{c}
e^{jx}=cos(x)+jsin(x) \\
cos(x)=cos(-x) \\
sin(x)=-sin(x) \\
e^{-jx}=cos(-x)+jsin(-x)=cos(x)-jsin(x) \\
\end{array}
\end{displaymath}
Possiamo quindi osservare che in $x(t)$ i termini $e^{\lambda_{3}t}$ e $e^{\lambda_{4}t}$ possono essere facilmente riscritti in modo che in essi compaiano esplicitamente i nostri modi:
\begin{displaymath}
\begin{array}{c}
e^{\lambda_{3}t}=e^{\sigma t-j\omega t}=e^{\sigma t}e^{-j\omega t}=e^{\sigma t}(cos(\omega t)-jsin(\omega t))=e^{\sigma t}cos(\omega t)-je^{\sigma t}sin(\omega t) \\
e^{\lambda_{4}t}=e^{\sigma t+j\omega t}=e^{\sigma t}e^{j\omega t}=e^{\sigma t}(cos(\omega t)+jsin(\omega t))=e^{\sigma t}cos(\omega t)+je^{\sigma t}sin(\omega t) \\
\end{array}
\end{displaymath}
E in generale, poiché noi abbiamo dei coefficienti a fattore, vale:
\begin{displaymath}
\begin{array}{c}
ke^{\lambda_{3}t}=ke^{\sigma t}cos(\omega t)-jke^{\sigma t}sin(\omega t) \\
ke^{\lambda_{4}t}=ke^{\sigma t}cos(\omega t)+jke^{\sigma t}sin(\omega t) \\
\end{array}
\end{displaymath}
Quindi per ogni coefficiente $k$ che moltiplica l'esponenziale complesso dobbiamo considerare una componente coseno (moltiplicata per 1) e una componente seno (moltiplicata per $-j$ o per $+j$). \\ \\
A questo punto esprimiamo l'evoluzione libera come richiesto, ovvero nella forma $x(t)=M(A,x_{0})m(t)$. Per farlo raggruppiamo i termini simili con opportuni colori e li posizioniamo nella matrice $M$ in modo che il prodotto con il vettore dei modi sia uguale alla soluzione libera:
\begin{displaymath}
M(A,x_{0})=\left(\begin{array}{cccc}
M_{11} & M_{12} & M_{13} & M_{14}\\
M_{21} & M_{22} & M_{23} & M_{24}\\
M_{31} & M_{32} & M_{33} & M_{34}\\
M_{41} & M_{42} & M_{43} & M_{44}\\
\end{array}\right)
\end{displaymath}
dove: \\
\begin{small}
\begin{displaymath}
\begin{array}{c}
M_{11}={\color{red}x_{0_{1}}+(-\frac{\beta\varepsilon}{\alpha\gamma})x_{0_{2}}}{\color{red}-\frac{1}{\alpha}x_{0_{3}}}{\color{red}+\frac{\beta}{\alpha\gamma}x_{0_{4}}}\\ \\
M_{12}={\color{green}\frac{1}{\alpha}(\frac{\beta(\varepsilon-\alpha)x_{0_{2}}-\beta x_{0_{4}}}{-\alpha^{2}+\alpha\varepsilon+\gamma}+x_{0_{3}})} \\ \\
M_{13}={\color{blue}\frac{2\beta(\sqrt{\Delta}+\varepsilon)}{\gamma(\sqrt{\Delta}-\varepsilon)(2\alpha+\sqrt{\Delta}-\varepsilon)}(\frac{2\gamma x_{0_{2}}-(\sqrt{\Delta}-\varepsilon)x_{0_{4}}}{2\sqrt{\Delta}})}+{\color{purple}\frac{2\beta(\sqrt{\Delta}-\varepsilon)}{\gamma(\sqrt{\Delta}+\varepsilon)(-2\alpha+\sqrt{\Delta}+\varepsilon)}(\frac{2\gamma x_{0_{2}}+(\varepsilon+\sqrt{\Delta})x_{0_{4}}}{2\sqrt{\Delta}})} \\ \\
M_{14}=-j{\color{blue}\frac{2\beta(\sqrt{\Delta}+\varepsilon)}{\gamma(\sqrt{\Delta}-\varepsilon)(2\alpha+\sqrt{\Delta}-\varepsilon)}(\frac{2\gamma x_{0_{2}}-(\sqrt{\Delta}-\varepsilon)x_{0_{4}}}{2\sqrt{\Delta}})}
+j{\color{purple}\frac{2\beta(\sqrt{\Delta}-\varepsilon)}{\gamma(\sqrt{\Delta}+\varepsilon)(-2\alpha+\sqrt{\Delta}+\varepsilon)}(\frac{2\gamma x_{0_{2}}+(\varepsilon+\sqrt{\Delta})x_{0_{4}}}{2\sqrt{\Delta}})} \\ \\
M_{21}=0 \\ \\
M_{22}=0 \\ \\
M_{23}=\textcolor{blue}{\frac{\sqrt{\Delta}+\varepsilon}{2\gamma}(\frac{2\gamma x_{0_{2}}-(\sqrt{\Delta}-\varepsilon)x_{0_{4}}}{2\sqrt{\Delta}})}+\textcolor{purple}{\frac{\sqrt{\Delta}-\varepsilon}{2\gamma}}{\color{purple}(\frac{2\gamma x_{0_{2}}+(\varepsilon+\sqrt{\Delta})x_{0_{4}}}{2\sqrt{\Delta}})} \\ \\
M_{24}=-j\textcolor{blue}{\frac{\sqrt{\Delta}+\varepsilon}{2\gamma}(\frac{2\gamma x_{0_{2}}-(\sqrt{\Delta}-\varepsilon)x_{0_{4}}}{2\sqrt{\Delta}})}+j\textcolor{purple}{\frac{\sqrt{\Delta}-\varepsilon}{2\gamma}}{\color{purple}(\frac{2\gamma x_{0_{2}}+(\varepsilon+\sqrt{\Delta})x_{0_{4}}}{2\sqrt{\Delta}})} \\ \\
M_{31}=0 \\ \\
M_{32}=\textcolor{green}{\frac{\beta(\varepsilon-\alpha)x_{0_{2}}-\beta x_{0_{4}}}{-\alpha^{2}+\alpha\varepsilon+\gamma}+x_{0_{3}}} \\ \\
M_{33}={\color{blue}\frac{\beta(\sqrt{\Delta}+\varepsilon)}{\gamma(2\alpha+\sqrt{\Delta}-\varepsilon)}(\frac{(\sqrt{\Delta}-\varepsilon)x_{0_{4}}-2\gamma x_{0_{2}}}{2\sqrt{\Delta}})}+\textcolor{purple}{\frac{\beta(\sqrt{\Delta}-\varepsilon)}{\gamma(-2\alpha+\sqrt{\Delta}+\varepsilon)}(\frac{2\gamma x_{0_{2}}+(\varepsilon+\sqrt{\Delta})x_{0_{4}}}{2\sqrt{\Delta}})} \\ \\
M_{34}=-j{\color{blue}\frac{\beta(\sqrt{\Delta}+\varepsilon)}{\gamma(2\alpha+\sqrt{\Delta}-\varepsilon)}(\frac{(\sqrt{\Delta}-\varepsilon)x_{0_{4}}-2\gamma x_{0_{2}}}{2\sqrt{\Delta}})}+j\textcolor{purple}{\frac{\beta(\sqrt{\Delta}-\varepsilon)}{\gamma(-2\alpha+\sqrt{\Delta}+\varepsilon)}(\frac{2\gamma x_{0_{2}}+(\varepsilon+\sqrt{\Delta})x_{0_{4}}}{2\sqrt{\Delta}})} \\ \\
M_{41}=0 \\ \\
M_{42}=0 \\ \\
M_{43}={\color{blue}\frac{(\sqrt{\Delta}-\varepsilon)x_{0_{4}}-2\gamma x_{0_{2}}}{2\sqrt{\Delta}}}+{\color{purple}\frac{2\gamma x_{0_{2}}+(\varepsilon+\sqrt{\Delta})x_{0_{4}}}{2\sqrt{\Delta}}} \\ \\
M_{44}=-j{\color{blue}\frac{(\sqrt{\Delta}-\varepsilon)x_{0_{4}}-2\gamma x_{0_{2}}}{2\sqrt{\Delta}}}+j{\color{purple}\frac{2\gamma x_{0_{2}}+(\varepsilon+\sqrt{\Delta})x_{0_{4}}}{2\sqrt{\Delta}}}\\ \\
\end{array}
\end{displaymath}
\end{small}
Anche in questo caso abbiamo verificato che la matrice risulta reale.
\end{solution}
\psimbC

\part[\afpt] Usando anche l'ipotesi di equilibrio (H2), discutere l'andamento asintotico dei modi e la stabilità del sistema.
\begin{solution} 
Utilizzando il vettore dei modi, composto dai modi ricavati precedentemente, ossia:
\begin{displaymath}
m(t)=\left(\begin{array}{c}
e^{\lambda_{1}t}\\
e^{\lambda_{2}t}\\
e^{\sigma t}\cos(\omega t)\\
e^{\sigma t}\sin(\omega t)
\end{array}\right)=\left(\begin{array}{c}
1\\
e^{\alpha t} \\
e^{\sigma t}\cos(\omega t)\\
e^{\sigma t}\sin(\omega t)
\end{array}\right)
\end{displaymath}
Passiamo ora ad analarizzare l'andamento (asintotico) di questi ultimi:
\begin{itemize}
\item Il primo elemento del vettore è un esponenziale semplice, del tipo $e^{\lambda t}$ con autovalore reale $\lambda_1=0$ , quindi è un modo costante e ha valore 1.
\item Il secondo elemento del vettore è un esponenziale semplice, del tipo $e^{\lambda t}$ con autovalore reale $\lambda_2=\alpha$; poichè $\alpha<0$, il modo è convergente a 0.
\item Il terzo e il quarto elemento del vettore sono funzioni oscillanti del tipo $e^{\sigma t}\cos{\omega t} / e^{\sigma t}\sin{\omega t}$ e sono quindi associati ad una coppia di autovalori complesi coniugati del tipo $\sigma \pm jw$; analizzando la parte reale caratterizzata da $\sigma$, notiamo che il suo valore è negativo, poichè è uguale a $\frac{\epsilon}{2}$ dove $\epsilon$ è dato dal rapporto $-\frac{kr}{j}$ il cui valore è <0, pertanto entrambi i modi sono convergenti a zero. 
\end{itemize}
Essendo che il sistema non possiede alcun modo divergente, ma possiede almeno un modo non convergente (ovvero $\lambda_1$ che è costante) possiamo dire che il sistema è marginalmente stabile.
\end{solution}
\psimbD

\part[\afpt] Calcolare la funzione di trasferimento $G(s)$ tra la variazione dell'ingresso e la variazione dell'uscita e determinare, se possibile, le posizioni dei rispettivi poli ed eventuali zeri
\begin{solution} 
% Si dice che la funzione di trasferimento è un invariante strutturale del sistema. Gli zeri di una funzione di trasferimento, ovvero le radici del polinomio a numeratore della f.d.t., sono legati alla presenza nella equazione differenziale che governa l’uscita delle derivate della funzione d’ingresso. \\
Esistono due metodi (tre contando anche la definizione distruttiva) per calcolare la funzione di trasferimento: applicare la trasformata di Laplace alle varie equazioni con le opportune sostituzioni oppure più semplicemente sfruttare la linearità della trasformata nelle matrici del modello linearizzato calcolate precedentemente. Possiamo quindi procedere come segue:
\begin{displaymath}
\left\{
\begin{array}{l}
sX(s) - x_0= AX(s) + BU(s)\\
Y(s)=CX(s) + DU(s)\\
\end{array}
\right.
\end{displaymath}
Possiamo porre la condizione iniziale $x_0=0$ e ricordare che $D=0$.
Adesso ricaviamo X(s) dalla prima equazione e la sostituiamo nella seconda in modo da ottenere Y(s):
\begin{displaymath}
\begin{array}{cc}
sX(s) - AX(s) = BU(s) \\
(sI-A)X(s) = BU(s) \\
(sI-A)^{-1}(sI-A)X(s) = (sI-A)^{-1}BU(s) \\
X(s) = (sI-A)^{-1}BU(s) \\
\end{array}
\end{displaymath}
Sostituendo: \\
\begin{displaymath}
Y(s)=CX(s)=C(sI-A)^{-1}BU(s)\\
\end{displaymath}
Da cui:
\begin{displaymath}
G(s)=\frac{Y(s)}{U(s)}=C(sI-A)^{-1}B
\end{displaymath}
Sostituendo le matrici ottenute in precedenza e dopo aver verificato che A è invertibile:
\begin{displaymath}
G(s)=\left(\begin{array}{cccc}
1 & -bcos(\bar{x_{2}}) & 0 & 0\end{array}\right)\left(\begin{array}{cccc}
s & 0 & -1 & 0\\
0 & s & 0 & -1\\
0 & -\beta & s-\alpha & 0\\
0 & -\gamma & 0 & s-\varepsilon
\end{array}\right)^{-1}\left(\begin{array}{c}
0\\
0\\
\frac{cos(\bar{x_{2}})}{M}\\
-\frac{b}{J}
\end{array}\right)
\end{displaymath}
Calcolando la inversa:
\begin{displaymath}
G(s)=C\left(\begin{array}{cccc}
\frac{1}{s} & \frac{\beta(s-\varepsilon)}{s(s-\alpha)(-\gamma+s^{2}-\varepsilon s)} & \frac{1}{s(s-\alpha)} & \frac{\beta}{s(s-\alpha)(-\gamma+s^{2}-\varepsilon s)}\\
0 & \frac{s-\varepsilon}{(-\gamma+s^{2}-\varepsilon s)} & 0 & \frac{1}{(-\gamma+s^{2}-\varepsilon s)}\\
0 & \frac{\beta(s-\varepsilon)}{(s-\alpha)(-\gamma+s^{2}-\varepsilon s)} & \frac{1}{(s-\alpha)} & \frac{\beta}{(s-\alpha)(-\gamma+s^{2}-\varepsilon s)}\\
0 & \frac{\gamma}{(-\gamma+s^{2}-\varepsilon s)} & 0 & \frac{s}{(-\gamma+s^{2}-\varepsilon s)}
\end{array}\right)B
\end{displaymath}
Svolgendo $C(sI-A)^{-1}$ otteniamo:
\begin{displaymath}
\begin{scriptsize}
\left(\begin{array}{cccc}
\frac{1}{s} & \frac{\beta(s-\varepsilon)}{s(s-\alpha)(-\gamma+s^{2}-\varepsilon s)}-bcos(\bar{x_{2}})\frac{s-\varepsilon}{(-\gamma+s^{2}-\varepsilon s)} & \frac{1}{s(s-\alpha)} & \frac{\beta}{s(s-\alpha)(-\gamma+s^{2}-\varepsilon s)}-bcos(\bar{x_{2}})\frac{1}{(-\gamma+s^{2}-\varepsilon s)}\end{array}\right)
\end{scriptsize}
\end{displaymath}
Moltiplicando il risultato ottenuto per la matrice B otteniamo:
\begin{displaymath}
G(s)=
\frac{1}{s(s-\alpha)}\frac{cos(\bar{x_{2}})}{M}+(\frac{\beta}{s(s-\alpha)(-\gamma+s^{2}-\varepsilon s)}-bcos(\bar{x_{2}})\frac{1}{(-\gamma+s^{2}-\varepsilon s)})(-\frac{b}{J})
\end{displaymath}
E semplificando:
\begin{scriptsize}
\begin{displaymath}
G(s)=\frac{(Jcos(\bar{x_{2}})+Mb^{2}cos(\bar{x_{2}}))s^{2}-(\varepsilon Jcos(\bar{x_{2}})+Mb^{2}cos(\bar{x_{2}})\alpha)s-(\gamma Jcos(\bar{x_{2}})+M\beta b)}{s(s-\alpha)(-\gamma+s^{2}-\varepsilon s)MJ}
\end{displaymath}
\end{scriptsize}
Calcolando gli zeri in forma simbolica:
\begin{displaymath}
\begin{array}{cc}
z_{1}=\frac{\varepsilon Jcos(\bar{x_{2}})+Mb^{2}cos(\bar{x_{2}})\alpha-\sqrt{(\varepsilon Jcos(\bar{x_{2}})+Mb^{2}cos(\bar{x_{2}})\alpha)^{2}+4(Jcos(\bar{x_{2}})+Mb^{2}cos(\bar{x_{2}}))(\gamma Jcos(\bar{x_{2}})+M\beta b)}}{2(Jcos(\bar{x_{2}})+Mb^{2}cos(\bar{x_{2}}))} \\
z_{2}=\frac{\varepsilon Jcos(\bar{x_{2}})+Mb^{2}cos(\bar{x_{2}})\alpha+\sqrt{(\varepsilon Jcos(\bar{x_{2}})+Mb^{2}cos(\bar{x_{2}})\alpha)^{2}+4(Jcos(\bar{x_{2}})+Mb^{2}cos(\bar{x_{2}}))(\gamma Jcos(\bar{x_{2}})+M\beta b)}}{2(Jcos(\bar{x_{2}})+Mb^{2}cos(\bar{x_{2}}))}
\end{array}
\end{displaymath}
E i poli:
\begin{displaymath}
\begin{array}{cc}
p_{1}=0 \\
p_{2}=\alpha \\
p_{3}=\frac{1}{2}(\varepsilon-\sqrt{\Delta}) \\
p_{4}=\frac{1}{2}(\varepsilon+\sqrt{\Delta})
\end{array}
\end{displaymath}
\end{solution}
\psimbE

\part[\afpt] Passando all'uso dei valori numerici indicati in Tabella~$2$, determinare valutare $G(s)$, discuterne la posizione e la molteplicità dei poli e degli eventuali zeri, e verificarne la stabilità.
\begin{solution} 
Sostituendo i valori numerici otteniamo:
\begin{displaymath}
\begin{array}{cc}
\alpha=-\frac{133*\frac{\sqrt{3}}{2}}{580}=\text{\ensuremath{-0.1986}} \\ \\
\beta=-\frac{-59.833*\frac{1}{2}}{580}+\frac{133*0*\frac{1}{2}}{580}-\frac{-143.20136*\frac{\sqrt{3}}{2}}{580}=0.2654 \\ \\
\gamma=-\frac{143*\frac{\sqrt{3}}{2}}{560}=\text{\ensuremath{-0.2211}} \\ \\
\varepsilon=-\frac{168}{560}=-0.3 \\
\end{array}
\end{displaymath}
\begin{small}
\begin{displaymath}
\begin{array}{cc}
G(s)=\frac{1}{s(s+0.1986)}\frac{\sqrt{3}}{2*580}+(\frac{0.2654}{s(s+0.1986)(0.2211+s^{2}+0.3s)}-1.2*\frac{\sqrt{3}}{2}\frac{1}{(0.2211+s^{2}+0.3s)})(-\frac{1.2}{560}) = \\ \\ 
=\frac{0.00372s^{2}+0.00089s-0.00023}{s\left(s+0.1986\right)\left(s^{2}+0.3s+0.2211\right)}
\approx\frac{0.0037201\left(s-0.1602\right)\left(s+0.3995\right)}{s\left(s+0.1986\right)\left(s^{2}+0.3s+0.2211\right)} 
\end{array}
\end{displaymath}
\end{small}
Calcoliamo gli zeri:
\begin{displaymath}
\begin{array}{cc}
z_1 = -0.3995 \\
z_2 = 0.1602  \\
\end{array}
\end{displaymath}
Calcoliamo i poli:
\begin{displaymath}
\begin{array}{cc}
p_1 = 0 \\ 
p_2 = -0.1986  \\
p_3 = -0.15 - 0.44564i \\
p_4 = -0.15 + 0.44564i \\
\end{array}
\end{displaymath}
Poichè $\Re(p_i)<0$, i poli sono stabili, giocoforza la funzione di trasferimento è stabile. \\
La molteplicità degli zeri e dei poli è 1.
\end{solution}
\psimbF

\end{parts}

%
%
%

\newpage
\section{Parte Numerica}

Al fine di progettare il controllore, si suggerisce di procedere secondo la metodologia vista durante il corso come segue:\\

\begin{parts}	

\part[\afpt] Tradurre i requisiti della Tabella~$3$ nelle corrispondenti specifiche statiche e dinamiche del sistema controllato, e indicare l'approssimazione $G_c(s)$ della f.d.t. del sistema complessivo risultante in anello chiuso.
\begin{solution} 
\textbf{Specifiche statiche:} \\
Segnali a gradino sono spesso usati per generare riferimenti
costanti. Per raggiungere l'inseguimento perfetto di un riferimento costante $r(t)$ per la variazione dell'uscita $y(t)$ è necessario imporre:
\[\lim_{t\to\infty} e(t)= \lim_{t\to\infty} (r(t)-y(t))=0 \]
Tuttavia tale condizione è verificata se $G(s)$ ha almeno un polo nell'origine, come nel nostro caso. \\ \\
\textbf{Specifiche dinamiche:} \\
Dobbiamo assicurare un inseguimento perfetto di un gradino con un tempo di risposta non superiore a 45s ed eventuali oscillazioni entro il 5\% dal valore di regime.
\\ Per farlo calcoliamo la $G_c(s)$. \\
Iniziamo utilizzando un'approssimazione ad un polo dominante ovvero: $$G_c(s)\approx \frac{1}{1+s\tau}$$
\begin{comment}
Ricaviamo lo smorzamento $\delta$ sfruttando la seguente relazione: 
\[S=e^{-\frac{\pi\delta}{\sqrt{1-\delta^2}}}\]
Svolgendo i seguenti passaggi: \\
\begin{displaymath}
\begin{array}{c}
\ln(S)=-\frac{\pi\delta}{\sqrt{1-\delta^2}} \\
(1-\delta^2)\ln^2(S)=\pi^2\delta^2 \\
\ln^2(S)-\delta^2\ln^2(S)=\pi^2\delta^2 \\
\ln^2(S)=\delta^2(\pi^2+\ln^2(S)) \\
\delta^2=\frac{\ln^2(S)}{\pi^2+\ln^2(S)} \\
\delta=\sqrt{\frac{\ln^2(S)}{\pi^2+\ln^2(S)}} \\
\end{array}
\end{displaymath}
E sostituendo i valori numerici otteniamo: \\ \\
$\delta=\sqrt{\frac{\ln^2(0.05)}{\pi^2+\ln^2(0.05)}}= 0.69010 \approx 0.7$ \\
\end{comment}
E' possibile ricavare il margine di fase $M_f$ richiesto sapendo che per i sistemi a un polo dominante la $\phi(t)=-90°$, ottenendo quindi $M_f=\pi+\phi (t) = 90°$.
\\
Ricaviamo adesso la pulsazione di taglio $\omega_T$ tramite la seguente relazione: \\ $$\omega_T \geq \frac{3}{T_a} = \frac{3}{45} = 0.0667 \, rad/s$$ \\
Si dovrà quindi ottenere un margine di fase $M_f \approx 90°$ e una frequenza di taglio $\omega_T \geq 0.0667 \, rad/s$.
\end{solution}
\pnumA

\part[\afpt] Se la f.d.t. $G(s)$ del modello in anello aperto non è stabile, progettare un anello di retroazione interno, caratterizzato da un controllore $C_i(s)$ e da una f.d.t. risultante $G_i(s)$; in tal caso, nell'elaborato devono essere forniti (1) la descrizione della procedura seguita per ottenere $C_i(s)$, e (2) le espressioni di $C_i(s)$ e $G_i(s)$ con le posizioni e molteplicità dei rispettivi poli e degli eventuali zeri.
\begin{solution} 
La f.d.t. $G(s)$ del modello in anello aperto è già stabile.
% Approndire cosa si doveva fare nel caso non fosse già stabile, a quanto pare montare qualcosa sulla catena
\end{solution}
\pnumB

\part[\afpt] Se invece $G(s)$ è già stabile o se si è ottenuta una $G_i(s)$ stabile dal punto precedente, verificare il soddisfacimento o meno delle specifiche statiche e dinamiche; a tal fine, riportare $G(s)$ o $G_i(s)$ in forma di Bode, indicando il guadagno di Bode, le costanti di tempo di eventuali singolarità reali, e i coefficienti di smorzamento e le pulsazioni naturali di eventuali singolarità coniugate.
\begin{solution} 
Riportiamo G(s) in forma di Bode: \\
\begin{displaymath}
G(s)\approx\frac{0.0037201\left(s-0.1602\right)\left(s+0.3995\right)}{s\left(s+0.1986\right)\left(s^{2}+0.3s+0.2211\right)}= \\
\end{displaymath}
Portando fuori un -0.1602:
\begin{displaymath}
=\frac{0.0037201(-0.1602)\left(1+s(-\frac{1}{0.1602})\right)\left(s+0.3995\right)}{s\left(s+0.1986\right)\left(s^{2}+0.3s+0.2211\right)}= \\
\end{displaymath}
Portando fuori un +0.3995:
\begin{displaymath}
=\frac{0.0037201(-0.1602)(+0.3995)\left(1+s(-\frac{1}{0.1602})\right)\left(1+s(\frac{1}{0.3995})\right)}{s\left(s+0.1986\right)\left(s^{2}+0.3s+0.2211\right)}= \\
\end{displaymath}
Portando fuori un +0.1986:
\begin{displaymath}
=\frac{0.0037201(-0.1602)(+0.3995)\left(1+s(-\frac{1}{0.1602})\right)\left(1+s(\frac{1}{0.3995})\right)}{s(+0.1986)\left(1+s(\frac{1}{0.1986})\right)\left(s^{2}+0.3s+0.2211\right)}= \\
\end{displaymath}
Per aggiustare il polinomio di secondo grado intanto facciamo spuntare +1 portando fuori un +0.2211:
\begin{displaymath}
=\frac{0.0037201(-0.1602)(+0.3995)\left(1+s(-\frac{1}{0.1602})\right)\left(1+s(\frac{1}{0.3995})\right)}{s(+0.1986)(+0.2211)\left(1+s(\frac{1}{0.1986})\right)\left(1+\frac{2}{1}\frac{0.3(2^{-1})}{0.2211}\frac{\sqrt{0.2211}}{\sqrt{0.2211}}s+\frac{1}{0.2211}s^2\right)}= \\
\end{displaymath}
Ricaviamo adesso $\delta=\frac{0.3(2^{-1}){\sqrt{0.2211}}}{0.2211}=0.31900$ e $\omega_n={\sqrt{0.2211}}=0.47021$.
\\ \\
Il guadagno di Bode, direttamente osservabile dalla forma di Bode, vale: \\
$$K=\frac{0.0037201(-0.1602)(+0.3995)}{(+0.1986)(+0.2211)}=-0.00542$$
Quindi la funzione di trasferimento in forma di Bode vale:
\begin{displaymath}
G(s)=\frac{-0.00542\left(1+s(-\frac{1}{0.1602})\right)\left(1+s(\frac{1}{0.3995})\right)}{s\left(1+s(\frac{1}{0.1986})\right)\left(1+\frac{2*0.319}{0.47021}s+\frac{1}{0.2211}s^2\right)}\\
\end{displaymath}
\\
Le costanti di tempo delle singolarità reali sono: \\
zeri: \\
$$\tau_{1}=-\frac{1}{z_{1}}=\frac{1}{0.3995}=2.50312$$
$$\tau_{2}=-\frac{1}{z_{2}}=-\frac{1}{0.1602}=-6.24219$$ 
poli: \\
$$\tau_{3}=-\frac{1}{p_{1}}=-\frac{1}{0}=-\infty$$ 
$$\tau_{4}=-\frac{1}{p_{2}}=\frac{1}{0.1986}=5.03524$$
I coefficienti di smorzamento e le pulsazioni naturali delle singolarità complesse coniugate sono già stati calcolati poco sopra. In alternativa ricordando che il coefficiente della parte reale è -0.15 e il coefficiente della parte immaginaria è 0.44564 si poteva ottenere lo stesso risultato calcolando: \\
$$\omega_n=\sqrt{\alpha^2+\beta^2}=\sqrt{(-0.15)^2+(0.44564)^2}=0.47020$$
$$\delta=-\frac{\alpha}{\omega_n}=-\frac{-0.15}{0.47020}=0.31901$$
\end{solution}
\pnumC

\part[\afpt] Progettare un controllore $C(s) = \frac{K_c}{s^t} C_0(s)$ che soddisfi i requisiti della Tabella~$3$ e ottenga la $G_c(s)$ desiderata; nel documento devono essere riportati (1) la descrizione dettagliata della procedura seguita, (2) le funzioni $C(s)$ e $G_c(s)$, nonché i (3) diagrammi di Bode nelle varie fasi della progettazione (ad esempio, quelli che mostrano le risposte armoniche della catena diretta $C(j\omega)\, G(j\omega)$ o $C(j\omega)\, G_i(j\omega)$), e del modello controllato complessivo $G_c(j\omega)$).
\begin{solution} 
(1) Descrizione: \\
Poiché la trasformata del gradino è 1/s è necessario avere un polo in zero per ottenere l'inseguimento perfetto.\\ 
Ma dal momento che tale polo è già presente in $G(s)$, non è necessario
aggiungerlo con il controllore e quindi t = 0. \\ \\
Essendo inizialmente il margine di fase $M_f=-93.1°$, come si può notare dalla \hyperref[fig:fig1]{Fig. 1}, risolviamo utilizzando un guadagno del controllore negativo $K_c=-1$ ottenendo un $M_f=86.9°$ e una frequenza di taglio $\omega_T = 0.00542 \, rad/s$, da come si può notare dalla \hyperref[fig:fig2]{Fig. 2}.\\ \\
Per aumentare la $\omega_T$, ovvero raggiungere il tempo di assestamento desiderato, procediamo con l'aumento del guadagno del controllore. \\
Andiamo a considerare quindi un guadagno pari a -11.7 ottenendo una frequenza di taglio $\omega_T=0.0671 \, rad/s$, da come si può notare dalla \hyperref[fig:fig3]{Fig. 3}.\\ \\
Infine aggiungiamo uno zero prima della frequenza di taglio e un polo in alta frequenza per aumentare il margine di fase, e inoltre ci possiamo permettere di abbassare di poco il guadagno per aumentare il margine di fase ancora un po', ottenendo il diagramma riportato in \hyperref[fig:fig4]{Fig. 4}.\\ \\
Per rendere un controllore fisicamente realizzabile è necessario che sia causale, ovvero che il numero di poli sia maggiore o uguale al numero di zeri, in caso contrario è necessario aggiungere poli in alta frequenza in modo da non perturbare il sistema. \\ \\
Essendo che nel nostro caso il numero dei poli è uguale al numero di zeri, il controllore è fisicamente realizzabile e quindi non è necessario aggiungere poli. \\
Si può notare dalla figura un margine di fase pari a 62.3° e una frequenza di taglio pari a 0.0672 rad/s.\\ \\
Il fatto che il margine di fase non sia uguale a 90° è dovuto all'approssimazione ad un polo dominante, che non è precisa essendoci poli vicini tra di loro.\\ \\
Possiamo notare che le specifiche richieste sono rispettate dalla risposta al gradino riportato nella figura \hyperref[fig:fig5]{Fig. 5}. \\

(2) Espressioni: \\
Otteniamo quindi il seguente controllore: $$C(s) = -11.4 \, \frac{1+\frac{s}{0.274}}{1+\frac{s}{0.902}}$$
Per ottenere la $G_c$ svolgiamo i seguenti passaggi:
$$G_c(s)=\frac{C(s)G(s)}{1+C(s)G(s)}$$
$$G_c(s)=\frac{-0.13961 \, (s+0.274)\,(s+0.3995)\,(s-0.1602)}{(s+0.9977)\,(s^2+0.2582s+0.01844)\,(s^2+0.1447s+0.1331)}$$ 
(3) Diagrammi di Bode: \\
(vedi paragrafo Figure e diagrammi)
\end{solution}
\pnumD

\part[\afpt] Realizzare una prima simulazione in Matlab/Simulink del sistema chiuso $G(s)$ in anello singolo con il controllore $C(s)$ o in doppio anello con i controllori $C_i(s)$ e $C(s)$; la simulazione deve mostrare il soddisfacimento dei requisiti della Tabella~$3$ rispetto ai segnali di riferimento o di disturbo indicati.
\begin{solution} 
Per svolgere la simulazione utilizziamo il seguente script matlab: 
\begin{scriptsize}
\lstinputlisting{support/runme.m}
\end{scriptsize}
\\ 
Dopo aver eseguito lo script selezioniamo il file Simulink simulazione.slx e eseguiamo la simulazione, che è possibile osservare nella figura \hyperref[fig:simulink]{Fig. 7} con uno stop time di 100s verificando che la risposta al gradino, riportata in \hyperref[fig:fig8]{Fig. 8}, corrisponde a quella in \hyperref[fig:fig6]{Fig. 6}
\end{solution}
\pnumE

\part[\afpt] Realizzare una seconda simulazione in Matlab/Simulink del modello del sistema non lineare, chiuso in anello singolo con il controllore $C(s)$ o in doppio anello con i controllori $C_i(s)$ e $C(s)$, operando le opportune traslazioni delle variabili in ingresso e in uscita dal sistema stesso; la simulazione deve mostrare il soddisfacimento dei requisiti della Tabella~$3$, anche se in modo approssimato, nei confronti dei soli segnali di riferimento.
\begin{solution} 
\input{2F.tex}
\end{solution}
\pnumF

\end{parts}

%
%
%

\newpage
\section{Parte Teorica}

\begin{parts}

\part[\afpt] \afTheoryQuestion
\begin{solution}
\input{3A.tex}
\end{solution}
\pteorA

\end{parts}

\newpage
\section*{Figure e diagrammi}
%%%%%%%%%%%%%%%%%%%%%%%%%%%%%%%%%%%%%%%%%%%%%%%%%%%%%%%%%%%%
%%%%%%%%%%%%%%%%%%% Add here all figures %%%%%%%%%%%%%%%%%%%
\begin{figure}[h]
\centering
\includegraphics[width=0.85\columnwidth]{img/fig1.eps}
\caption{Diagrammi di Bode di $C(s) \, G(s)$ con $C(s)=1$.}
\label{fig:fig1}
\end{figure}

\begin{figure}[h]
\centering
\includegraphics[width=0.85\columnwidth]{img/fig2.eps}
\caption{Diagrammi di Bode di $C(s) \, G(s)$ con $C(s)=-1$.}
\label{fig:fig2}
\end{figure}

\newpage

\begin{figure}[h]
\centering
\includegraphics[width=0.85\columnwidth]{img/fig3.eps}
\caption{Diagrammi di Bode di $C(s) \, G(s)$ con $C(s)=-11.7$.}
\label{fig:fig3}
\end{figure}

\begin{figure}[h]
\centering
\includegraphics[width=0.85\columnwidth]{img/fig4.eps}
\caption{Diagrammi di Bode di $C(s) \, G(s)$ con $C(s)=-11.4 \, \frac{1+3.6s}{1+1.1s}$.}
\label{fig:fig4}
\end{figure}

\newpage

\begin{figure}[h]
\centering
\includegraphics[width=0.85\columnwidth]{img/fig5.eps}
\caption{Risposta al gradino con caratteristiche indicate}
\label{fig:fig5}
\end{figure}

\begin{figure}[h]
\centering
\includegraphics[width=0.85\columnwidth]{img/fig6.eps}
\caption{Risposta al gradino}
\label{fig:fig6}
\end{figure}

\newpage

\begin{figure}[h]
\centering
\includegraphics[width=1\columnwidth]{img/simulazione.pdf}
\caption{Schema Simulink che consente la simulazione e la verifica del soddisfacimento delle specifiche.}
\label{fig:simulink}
\end{figure}

\newpage

\begin{figure}[h]
\centering
\includegraphics[width=1\columnwidth]{img/fig8.eps}
\caption{Risposta al gradino con Stop Time = 100s.}
\label{fig:fig8}
\end{figure}
%%%%%%%%%%%%%%%%%%%%%%%%%%%%%%%%%%%%%%%%%%%%%%%%%%%%%%%%%%%%
%%%%%%%%%%%%%%%%%%%%%%%%%%%%%%%%%%%%%%%%%%%%%%%%%%%%%%%%%%%%


\finalComments

\vfill \medskip\hfill\textbf{Fine dello svolgimento}

\end{document}

